% ============================================================
%  preamble.tex -- packages, colours, boxes and macros
%  (study-book skill template; fill the @@...@@ placeholders)
% ============================================================
\usepackage[a4paper,inner=24mm,outer=34mm,top=24mm,bottom=26mm,
            marginparwidth=24mm,marginparsep=3mm]{geometry}
\usepackage[T1]{fontenc}
\usepackage[utf8]{inputenc}
\usepackage[@@BABEL_LANG@@]{babel}
\usepackage{lmodern}
\usepackage[scaled=0.92]{helvet}
\usepackage{microtype}
\usepackage{amsmath,amssymb}
\usepackage{textcomp}
\usepackage[version=4]{mhchem}
\usepackage{siunitx}
\sisetup{range-phrase=\text{--},range-units=single,per-mode=symbol,
         separate-uncertainty=true,group-minimum-digits=4}
\usepackage{booktabs,tabularx,longtable,array,multirow}
\usepackage{graphicx}
\usepackage{adjustbox}
\graphicspath{{figures/}}
\usepackage{float}
\usepackage{caption}
\captionsetup{font=small,labelfont=bf,format=hang}
\usepackage{subcaption}
\usepackage[dvipsnames,table]{xcolor}
\usepackage{enumitem}
\setlist{itemsep=1pt,topsep=3pt,parsep=0pt}
\usepackage{fancyhdr}
\usepackage{titlesec}
\usepackage{marginnote}
\usepackage{listings}
\usepackage{tikz}
\usetikzlibrary{arrows.meta,positioning,shapes.geometric,shapes.symbols,
  shapes.misc,calc,decorations.pathreplacing,fit,backgrounds,patterns}
\usepackage{pgfplots}
\pgfplotsset{compat=1.18}
\usepackage[most]{tcolorbox}
\tcbuselibrary{listings}
\usepackage[hidelinks]{hyperref}
\hypersetup{colorlinks=true,linkcolor=NavyBlue,urlcolor=NavyBlue,citecolor=NavyBlue,
  pdftitle={@@TITLE@@},pdfauthor={@@AUTHOR@@}}
\usepackage[nameinlink,capitalise]{cleveref}

% ---------------- course-specific macros / units ----------------
% (the main agent adds \DeclareSIUnit, \newcommand etc. from glossary.md here)
@@COURSE_MACROS@@

% ---------------- colours ----------------
\definecolor{brand}{HTML}{1F4E79}
\definecolor{critical}{HTML}{B3261E}
\definecolor{lecq}{HTML}{C77700}
\definecolor{beyond}{HTML}{5F6B73}
\definecolor{pit}{HTML}{6A1B9A}
\definecolor{flow}{HTML}{2E75B6}
\definecolor{flowlight}{HTML}{DCE9F5}
\definecolor{codebg}{HTML}{F6F8FA}
\definecolor{unsure}{HTML}{E65100}

% ---------------- headings ----------------
\titleformat{\chapter}[display]
  {\normalfont\sffamily\huge\bfseries\color{brand}\raggedright}
  {\large\chaptertitlename\ \thechapter}{8pt}{\Huge}
\titlespacing*{\chapter}{0pt}{-10pt}{24pt}
\titleformat{\part}[display]{\centering\normalfont\sffamily\Huge\bfseries\color{brand}}
  {\LARGE\partname\ \thepart}{18pt}{}
\titleformat{\section}{\raggedright\normalfont\sffamily\Large\bfseries\color{brand}}{\thesection}{0.8em}{}
\titleformat{\subsection}{\raggedright\normalfont\sffamily\large\bfseries\color{brand!85!black}}{\thesubsection}{0.8em}{}
\titleformat{\subsubsection}{\normalfont\sffamily\normalsize\bfseries}{\thesubsubsection}{0.8em}{}
\setcounter{secnumdepth}{2}
\setcounter{tocdepth}{1}

\pagestyle{fancy}
\fancyhf{}
\fancyhead[LE]{\sffamily\small\leftmark}
\fancyhead[RO]{\sffamily\small\rightmark}
\fancyfoot[LE,RO]{\sffamily\small\thepage}
\fancyfoot[C]{\sffamily\scriptsize\color{gray}@@FOOTER@@}
\renewcommand{\headrulewidth}{0.4pt}
\setlength{\headheight}{16pt}

\setlength{\parindent}{0pt}
\setlength{\emergencystretch}{3em}
\raggedbottom
\setlength{\parskip}{4pt plus 1pt}
\renewcommand{\arraystretch}{1.18}

% ---------------- boxes ----------------
\tcbset{studybox/.style={enhanced,breakable,boxrule=0.6pt,arc=2pt,
  left=6pt,right=6pt,top=4pt,bottom=4pt,fonttitle=\sffamily\bfseries\small,
  before skip=8pt,after skip=8pt}}

\newtcolorbox{examcritical}[1][]{studybox,colback=critical!5,colframe=critical,
  title={\textbf{!}\ \ Exam-critical#1}}
\newtcolorbox{beyondslides}[1][]{studybox,colback=beyond!6,colframe=beyond,
  title={Beyond the sources#1},fontupper=\small}
\newtcolorbox{pitfall}[1][]{studybox,colback=pit!5,colframe=pit,
  title={Pitfall / common confusion#1}}
\newtcolorbox{keyidea}[1][]{studybox,colback=brand!5,colframe=brand,
  title={Key idea#1}}

% Inline citation of a source page: \src{TAG}{23}, \src{TAG}{23--25}, \src{TAG}{3,7}
\newcommand{\src}[2]{\unskip\nobreak\hspace{0.4em}{\sffamily\scriptsize\color{gray}[#1~#2]}}

% Lecturer (highlighted) question with worked answer.
% \begin{lecturerq}{<key>:label}{Question text}{TAG}{page} ... answer ... \end{lecturerq}
\newcounter{lecq}
\newenvironment{lecturerq}[4]{%
  \refstepcounter{lecq}\label{lq:#1}%
  \begin{tcolorbox}[studybox,colback=lecq!6,colframe=lecq,
    title={Lecturer question \thelecq\hfill{\normalfont\footnotesize [#3~#4]}}]
  \textbf{\itshape ``#2''}\par\smallskip
  \textbf{Model answer.}\ }{\end{tcolorbox}}

% A value that could not be confirmed (second read disagreed / unreadable source).
% \unsure{1.5 kV}  -> printed in orange with a warning mark; listed in the corrections appendix
\newcommand{\unsure}[1]{{\color{unsure}#1\textsuperscript{\,\textbf{?!}}}}

% Learning objectives at chapter start
\newenvironment{objectives}{%
  \begin{tcolorbox}[studybox,colback=white,colframe=brand!50,
     title={After this chapter you can\ldots}]\begin{itemize}[leftmargin=1.2em]}%
  {\end{itemize}\end{tcolorbox}}

% Take-home, self-test and answers
\newcommand{\takehome}{\section*{Take-home messages}}
\newcommand{\selftest}{\section*{Self-test}}
\newenvironment{answers}{\subsection*{Answers}\small\begin{enumerate}[label=\arabic*.]}{\end{enumerate}}

% ---------------- code (notebooks) ----------------
\lstdefinestyle{book}{basicstyle=\ttfamily\footnotesize,breaklines=true,columns=fullflexible,
  keepspaces=true,showstringspaces=false,commentstyle=\color{gray}\itshape,
  keywordstyle=\color{brand}\bfseries,stringstyle=\color{critical!80!black},
  upquote=true,literate={-}{-}1}
% \begin{codebox}[Python]{Title} ... \end{codebox}
\newtcblisting{codebox}[2][Python]{studybox,colback=codebg,colframe=flow!60,
  listing only,listing options={style=book,language=#1},title={\ttfamily #2}}

% ---------------- workflow chevrons ----------------
% \workflow{Step 1,Step 2,Step 3}
\tikzset{wfstep/.style={signal,signal to=east,signal from=west,draw=flow,
   fill=flowlight,text=brand,font=\sffamily\scriptsize\bfseries,
   minimum height=9mm,inner xsep=5pt,align=center,text width=#1},
   wfstep/.default=16mm}
\newcommand{\workflow}[2][]{%
  \foreach \s [count=\n] in {#2}{\xdef\wfn{\n}}%
  \pgfmathsetlengthmacro\wfw{min(24mm,(\linewidth-\wfn*10mm)/\wfn)}%
  \begin{center}\begin{adjustbox}{max width=\linewidth}\begin{tikzpicture}[node distance=1.2mm]
  \foreach \s [count=\i] in {#2}{%
    \ifnum\i=1 \node[wfstep=\wfw] (w\i) {\s};
    \else \pgfmathtruncatemacro\prev{\i-1}\node[wfstep=\wfw,right=of w\prev] (w\i) {\s};\fi}
  \end{tikzpicture}\end{adjustbox}\end{center}}

% ---------------- overview map ----------------
% solid arrow = dependency stated in the sources; dashed = inferred by the editor
\tikzset{
  ovpart/.style={draw=brand,fill=brand!8,rounded corners=3pt,inner sep=6pt},
  ovchap/.style={draw=flow,fill=flowlight,rounded corners=2pt,font=\sffamily\scriptsize\bfseries,
     text=brand,align=center,text width=30mm,minimum height=9mm},
  ovdep/.style={-{Stealth[length=2mm]},thick,brand},
  ovinf/.style={-{Stealth[length=2mm]},thick,brand!60,dashed},
}

% ---------------- misc helpers ----------------
\newcommand{\ie}{i.e.\ }
\newcommand{\eg}{e.g.\ }
\crefname{lecq}{lecturer question}{lecturer questions}
\Crefname{lecq}{Lecturer question}{Lecturer questions}

% Blank verso pages (two-sided layout) carry no header/footer
\makeatletter
\renewcommand{\cleardoublepage}{\clearpage\if@twoside\ifodd\c@page\else
  \hbox{}\thispagestyle{empty}\newpage\fi\fi}
\makeatother
